% =========================================================
\section{引言}
\label{sec:introduction}
% =========================================================

储能电站的运维（O\&M）知识具有天然的动态性和过程性。处理单个异常往往需要联合查阅设备状态、告警与巡检记录、历史工单、操作规程以及处置后的验证结果。与静态知识问答不同，系统在判断哪些证据与查询相关之外，还必须回答两个更严格的问题：这些证据在查询时刻是否已经可用，以及这些证据能否共同构成支持当前决策的完整过程。

检索增强生成（Retrieval-Augmented Generation, RAG）将语言模型与外部检索结合，用于知识密集型生成~\cite{lewis2020rag}。在检索端，BM25 仍是较强的词法基线，而 DPR、ColBERT 等神经检索器分别通过稠密表示和后期交互提升语义匹配能力~\cite{karpukhin2020dpr,khattab2020colbert}；BEIR 进一步表明，不同检索范式在跨领域泛化能力上存在明显差异~\cite{thakur2021beir}。在生成端，FiD 表明生成器能够有效聚合多段检索文本~\cite{izacard2021fid}，Self-RAG 则进一步研究了对检索、生成与自反思过程的联合控制~\cite{asai2024selfrag}。

然而，传统 RAG 通常将候选文档视为彼此独立的相关性单元，这一假设并不适用于持续更新的运维记录。首先，运维证据具有双重时间属性：事件发生时间与该记录对系统可用的时间并不一定一致，晚到记录不能被早先的查询追溯使用。其次，运维决策通常依赖互补证据：状态观测触发巡检，巡检支持诊断，诊断导出处置，而处置后的复查用于验证前述决策。因此，即使单条文档得分很高，也不能保证 Top-$k$ 证据集合覆盖完整过程。

已有工作从不同角度处理了这一问题的相关方面。时间感知 RAG 对时间约束或跨时间区间的信息变化进行建模~\cite{zhang2025mrag,lau2025tarag}；多跳检索根据先前已检索的信息来决定后续检索步骤~\cite{xiong2021mdr,trivedi2023ircot}；RAPTOR、GraphRAG 和 HippoRAG 则利用层次结构或图结构组织跨文档信息~\cite{sarthi2024raptor,edge2024graphrag,gutierrez2024hipporag}。这些工作分别从时间过滤、迭代检索与跨文档组织等方面提供了相关解决思路。

本文提出 TEF-RAG（Temporal Evidence Flow Retrieval-Augmented Generation）。其核心思想是在查询时刻可见的证据之间构建由当前查询条件化的有向语义关系，并将检索从独立节点排序重新表述为证据集合排序。系统首先利用确定性的资产、时间与规程约束剔除不可用记录；随后学习哪些候选证据对值得进行关系推理，并据此构建时序证据图；最后对大小受限的候选集合进行非线性成对排序，使模型学习当前查询下哪个候选集合更符合标注的完整性目标。

TEF-RAG 与一般图检索的区别在于，其图结构并非由静态实体关系预先定义，而是基于当前查询下实际可见的运维证据动态形成。它与时间感知排序的区别在于，时间约束仅用于定义证据是否合法，最终目标仍是恢复一个角色互补且关系一致的证据集合。它也不同于逐步式多跳检索，因为目标结构可以包含分支、更新、替代以及保留的不确定性，而不局限于单一路径。

在本文考虑的具身运维场景中，检索不是终点。选出的证据流会被传递给一个共享生成器，用于生成标准化工单和具有依赖关系的动作计划，从而为下游规划器、约束检查器和机器人执行模块提供基于证据的接口。本文重点研究该流程中的证据检索和结构化规划上下文，并不声称已经解决闭环机器人控制或现场安全认证问题。

本文的主要贡献有三点：
\begin{enumerate}
    \item 我们面向动态运维知识提出时序证据流（Temporal Evidence Flow）检索，将时间可见性与证据集合完整性统一到同一检索框架中。
    \item 我们提出一种结合学习式证据对提议、查询条件化关系图与非线性集合排序的方法，将检索目标从独立文档相关性转向集合级过程效用。
    \item 我们构建了一个受控的时序运维基准，并在共享生成器下同时使用检索指标和结构化生成指标，评估证据流质量是否反映到下游工单与动作计划中。
\end{enumerate}

% =========================================================
\subsection{相关工作}
\label{sec:related_work}
% =========================================================

\textbf{检索增强生成与神经检索。}
经典 RAG 将外部检索器与生成器结合~\cite{lewis2020rag}；DPR 和 ColBERT 分别代表稠密双编码器检索与后期交互检索~\cite{karpukhin2020dpr,khattab2020colbert}。FiD 从生成端展示了多证据聚合能力~\cite{izacard2021fid}，而 BEIR 系统比较了稀疏检索、稠密检索与重排序范式的泛化能力~\cite{thakur2021beir}。DPR 和 ColBERT 以单段文本为单位进行排序，FiD 则在生成阶段聚合多段已检索文本。相比之下，我们的检索目标直接对一个证据集合及其时间与语义关系赋予效用。

\textbf{时间感知与多跳检索。}
MRAG 和 TA-RAG 将时间条件显式纳入检索过程~\cite{zhang2025mrag,lau2025tarag}；MDR 和 IRCoT 通过逐步检索或检索与推理交替的方式寻找多条证据~\cite{xiong2021mdr,trivedi2023ircot}。TEF-RAG 同样强调互补证据，但进一步区分事件时间与可用时间，并允许证据结构包含更新、替代、验证以及不确定性传播。

\textbf{结构化与图式 RAG。}
RAPTOR、GraphRAG 和 HippoRAG 利用层次化或图结构组织跨文档信息~\cite{sarthi2024raptor,edge2024graphrag,gutierrez2024hipporag}。与通常从文档层次或实体关系构图的方法不同，TEF-RAG 的边由当前查询与当前可见的运维记录共同决定，其语义直接对应运维证据之间的支持、更新、限定和验证等关系。

\textbf{具身运维与电池智能。}
RAG 也开始用于现场机器人和专业设备操作。SayComply 研究现场机器人对操作合规知识的检索，RAGLRO 则将检索增强语言模型应用于机器人操作任务~\cite{ginting2025saycomply,wang2026raglro}。在电池领域，LiBrain 探索了将时间序列信息、检索知识与语言模型结合用于电池诊断和管理~\cite{li2026librain}。这些研究将检索到的操作知识或时间序列信息连接到机器人或电池管理任务。本文聚焦其前置的证据选择环节，使下游工单和动作计划所使用的证据在查询时刻可见，并覆盖任务所需的过程角色。

\textbf{学习排序。}
TEF-RAG 的最终选择阶段是查询级学习排序问题。RankNet 从成对偏好中学习相对排序函数~\cite{burges2005ranknet}，ANCE 等工作则表明，难负样本对学习有效的检索边界十分重要~\cite{xiong2021ance}。我们将这一思想应用到证据集合而非单条文档：训练目标在同一查询下直接比较完整集合与高分但不完整的难负集合。
